\section{Problem Context}
% Contexto institucional >>> limitações espaciais >>> assimetria temporal >>> desajuste dos incentivos.

% Literatura aplicada atual

% ---> Mecanismos atuais de governança ambiental (regulação, precificação, mercados de carbono, etc.) são insuficientes para lidar com a complexidade e a dinâmica dos sistemas socioecológicos.

% ---> Dados e evidências mostram que, apesar dos esforços, as tendências de degradação ambiental e perda de biodiversidade continuam, indicando falhas estruturais nos modelos de governança atuais.

% ---> Citar esforços internacionais (Acordo de Paris, Objetivos de Desenvolvimento Sustentável) e nacionais (políticas ambientais, leis de proteção) que não têm conseguido reverter as tendências negativas.

% ---> Usar dados de desmatamento, perda de biodiversidade, emissões de gases de efeito estufa, Brasil e mundo, para ilustrar a insuficiência dos mecanismos atuais.

% --> Concluir que é necessário repensar os fundamentos teóricos e práticos da governança ambiental, buscando abordagens que integrem melhor a natureza como parte dos sistemas econômicos e sociais, e que sejam mais adaptativas, inclusivas e eficazes.

In response to the ongoing environmental degradation, over the past decades, a range of institutional efforts has been undertaken to curb environmental disasters such as global warming, tropical deforestation, and the indiscriminate or illegal exploitation of water and soil resources. International agreements such as the Kyoto Protocol (1997), the Paris Agreement (2015), and the Convention on Biological Diversity (1992), as well as regional initiatives and national policies, have established frameworks for emission reduction, biodiversity conservation, and sustainable resource management. Not necessarily considering human self-interest over ecological needs, these instruments often rely on legally binding commitments, emission trading systems, reforestation programs, and technological cooperation mechanisms \cite{easterly2013tyranny}.

In a related vein, legal and institutional instruments have been established to create a system of checks and balances capable of inhibiting irrational environmental exploitation in the short and medium term, while fostering ecological awareness in the long term. For example, the Brazilian Forest Code (Lei nº 12.651/2012) stands out as a key legal framework regulating land use and native vegetation preservation in rural properties within the country's borders \cite{brasil2012florestal}, while non-governmental organizations (NGOs) such as Greenpeace, the World Wide Fund for Nature (WWF), the Rainforest Foundation, Imazon (Instituto do Homem e Meio Ambiente da Amazônia), among many others, play a crucial role in monitoring compliance, raising public awareness, and pressuring both governments and corporations to adopt more sustainable practices throughout the world \cite{lomborg2001skeptical}. 

Despite their relevance, such efforts often reflect a moralistic or idealized form of environmentalism. While rhetorically aimed at fostering mutualistic relationships between humans and nature, they often result in a regulated form of extractivism. Legal frameworks such as the Forest Code impose rigid constraints that, while intended to protect native ecosystems, fail to align with the economic rationale of landowners, who often perceive them not as incentives for sustainable coexistence, but as punitive restrictions that prevent them from fully utilizing a productive factor — land — under their control \cite{agrawal2005environmentality}. This maladjustment not only generates resistance and non-compliance but also fuels a broader sense of injustice among rural producers, who feel disproportionately burdened by conservation obligations without corresponding compensatory mechanisms \cite{borner2016mixing}.

Theoretical and applied economics, in turn, has addressed environmental issues with the common understanding that markets fail to internalize ecological costs. Without some form of coercive intervention — usually by governments — such failures tend to generate vicious long-term outcomes. For instance, \citeonline{pigou2010liberty} have proposed tariffs (Pigouvians taxes) as a way to internalize negative environmental externalities alongside its neoclassical model. In his view, \citeonline{hardin1968tragedy} argued that, in the absence of private property or state regulation, individual self-interest leads to what he called the "tragedy of the commons", a situation in which natural resources are completely exhausted as a result of individuals aiming to maximize their own utility, while the common good is entirely neglected. 

Since the mid-1980s, the so-called Green Economy has advanced the argument that economic development is only viable if improvements in human well-being are balanced with reductions in environmental risks and ecological scarcity \cite{jacobs1991green}. According to this school of thought, \textbf{the economy must stop treating the environment as a substitute good and begin to value it as a complementary asset}. This implies ensuring that each generation inherits an ecosystem that is at least equivalent in quality and function, recognizing that many ecological characteristics are irreplaceable. Such an approach is framed not only by the risks of irreversibility and uncertainty but also by the moral imperative of intergenerational equity \cite{pearce1989blueprint}.

As we can see, throughout history, a persistent misalignment has existed between collective supply and private demand in the use of natural resources, which leads to misconduct on the part of certain private agents who do not consider it fair to internalize diffuse social costs in individual operations, all of this intertwined in a non-standardized and inefficient institutional-legal scheme, operated by States and monitored by the third sector. Despite advances in theory and policy, a deeper problem persists: the economic system — including markets, institutions, and even economic theory — struggles to adequately conceptualize and resolve the trade-off between development and sustainability. 

Apparently, an endless back-and-forth dilemma is presented as an unsolvable tension between collective morality and individual utilitarianism. In the end, nature seems to be on a path of degradation that initially affects vulnerable regions, but eventually reaches everyone, as environmental collapse leads to widespread economic consequences such as global warming, water scarcity, and the increasing threat to agricultural and ecological food systems. This apparent antagonism may in fact be rooted in theoretical and institutional misalignments, rather than from an inherent incompatibility between individual incentives and ecological responsibility, a conjecture explored throughout this study.